% Source-derived chapter 08: Hitchin-type moduli spaces
% Original source: higher-siegel-weil-unitary-nonsingular-terms
\chapter{Hitchin-type moduli spaces}
\label{higher-siegel-weil-unitary-nonsingular-terms-chapter-08}
\source{higher-siegel-weil-unitary-nonsingular-terms}

\begin{remark}[Source section]
\label{higher-siegel-weil-unitary-nonsingular-terms-chapter-08-source}
\source{higher-siegel-weil-unitary-nonsingular-terms}
\source{higher-siegel-weil-unitary-nonsingular-terms}
This chapter reproduces the corresponding section of the retrieved
LaTeX source, including its definitions, statements, examples, and
proofs. It has not yet been translated into Lean.
\notready
\end{remark}

% The source section title was: Hitchin-type moduli spaces
\label{sec: hitchin spaces}
\source{higher-siegel-weil-unitary-nonsingular-terms}

In this section we introduce certain ``Hitchin-type moduli stacks'' which will help to analyze the special cycles. In particular, we will be able to use these to give an alternative construction of the cycle classes associated to special divisors, that is manifestly independent of auxiliary choices. 

\subsection{Hitchin stacks}\label{ssec: Hitchin stacks}
\source{higher-siegel-weil-unitary-nonsingular-terms}
Until \S\ref{ss:vb to tor}, we fix an arbitrary positive integer $m$.

\begin{defn}
\source{higher-siegel-weil-unitary-nonsingular-terms}
\notready The \emph{Hitchin stack} $\Cal{M}^{\rm all}(m,n)$ (sometimes denoted $\Cal{M}^{\rm all}$ when $m,n$ are understood) has $S$-points the groupoid consisting of the following data. 
\begin{itemize}
\item $\Cal{E}$ a rank $m$ vector bundle on $X' \times S$.
\item $\Cal{F}$ a rank $n$ vector bundle on $X' \times S$, equipped with a Hermitian map $h \co \Cal{F} \xrightarrow{\sim} \sigma^* \Cal{F}^{\vee}$. 
\item A map of underlying coherent sheaves $t \co \Cal{E} \rightarrow \Cal{F}$ over $X'\times S$.\end{itemize}
We define  $\cM(m,n)\subset\cM^{\rm all}$ (sometimes denoted $\Cal{M}$ when $m,n$ are understood) to be the open substack where the map $t$ base changes to an injective map on $X'_{\ol{s}}$ for each geometric point $\ol{s} \rightarrow S$. 
\end{defn}
Let us emphasize that \emph{both} $\Cal{E}$ and $(\Cal{F},h)$ are varying in this definition. We will usually suppress the dependence on $m,n$ from the notation. 


\subsection{Hitchin base} \label{ss:Hit base}
\source{higher-siegel-weil-unitary-nonsingular-terms}
%\WZ{Define $\cA_\cE$ in \S7, $\cA^{all}(m,n)$.}
\begin{defn}
\source{higher-siegel-weil-unitary-nonsingular-terms}
\notready We define the following two versions of the \emph{Hitchin base}. 
\begin{enumerate}
\item $\cA^{\mrm{all}}(m)$ (sometimes denoted $\cA^{\mrm{\all}}$ when $m$ is understood) to be the stack whose $S$-points is the groupoid of the following data: 
\begin{itemize}
\item $\Cal{E}$ a rank $m$ vector bundle on $X'\times S$;
\item $a \co \Cal{E} \rightarrow \sigma^* \Cal{E}^{\vee}$ is a map of coherent sheaves on $X'\times S$ such that $\sigma(a^{\vee}) = a$.
\end{itemize}

\item We define $\cA \subset \cA^{\mrm{all}}$ to be the open substack where $a\co \Cal{E} \rightarrow \sigma^* \Cal{E}^{\vee}$ is injective after base change to $X'_{\ol{s}}$ for every geometric point $\ol{s} \rightarrow S$.
\end{enumerate}
\end{defn}

\begin{defn}
\source{higher-siegel-weil-unitary-nonsingular-terms}
\notready For integers $1 \leq m \leq n$, we define the \emph{Hitchin fibration} for $\cM^{\rm all} = \cM^{\rm all}(m,n)$ to be the map $f \co \cM^{\rm all} \rightarrow \Cal{A}^{\mrm{all}}$ sending $(\Cal{E}, (\Cal{F},h), t)$ to
the composition 
\[
a\co \Cal{E} \xrightarrow{t} \Cal{F} \xrightarrow{h} \sigma^* \Cal{F}^{\vee} \xrightarrow{\sigma^* t^{\vee}} \sigma^* \Cal{E}^{\vee}.
\]	
\end{defn}



\begin{remark}
\source{higher-siegel-weil-unitary-nonsingular-terms}
\notready\label{r:M to A}In general the Hitchin fibration does not send $\cM(m,n)$ to $\cA(m)$ even when $m\le n$. However, in the special case $m=n$,  the Hitchin map does send $\cM(n,n)$ to $\cA(n)$ because when $t:\cE\to \cF$ is generically injective, it is generically an isomorphism for rank reasons, hence the induced Hermitian map $a$ on $\cE$ is generically non-degenerate.
\source{higher-siegel-weil-unitary-nonsingular-terms}
\end{remark}


%We have $\Cal{M}_{\GL(m), U(n)} = \coprod_d \Cal{M}_{\GL(m),U(n)}^d$ with $ \Cal{M}_{\GL(m),U(n)}^d$ being the locus where the map $\Cal{E} \rightarrow \Cal{F}$ has colength $d$. \tony{perhaps this remark belongs somewhere later} 


\subsection{Hitchin shtukas}\label{ssec: Hitchin shtukas}
\source{higher-siegel-weil-unitary-nonsingular-terms}
We now discuss a notion of shtukas for Hitchin stacks. Throughout, $\cM = \cM(m,n)$. 

\begin{defn}[Hecke stacks for Hitchin spaces]
\source{higher-siegel-weil-unitary-nonsingular-terms}
\notready\label{defn: hitchin hecke}
\source{higher-siegel-weil-unitary-nonsingular-terms}
For $r \geq 0$, we define $\Hk_{\Cal{M}^{\all}}^{r}$ to be the stack whose $S$-points are given by the groupoid of the following data: 	
\begin{enumerate}
\item $(\{x'_i\}_{1\le i\le r},\{(\cF_{i},h_{i})\}_{0\le i\le r})\in \Hk^{r}_{U(n)}(S)$. 
\item A vector bundle $\Cal{E}$ of rank $m$ on $X' \times S$. 
\item Maps $t_i \co \Cal{E} \rightarrow \Cal{F}_i$ fitting into the commutative diagram
\end{enumerate}
\[
\begin{tikzcd}
\Cal{E} \ar[r, equals] \ar[d, "t_0"] & \Cal{E} \ar[r, equals]  \ar[d, "t_1"] & \ldots \ar[r, equals] & \Cal{E} \ar[d, "t_r"]  \\
\Cal{F}_{0}  \ar[r, dashed] & \Cal{F}_1 \ar[r, dashed] & \ldots \ar[r, dashed] &  \Cal{F}_r
\end{tikzcd}
\]
We define the open substack $\Hk_{\Cal{M}}^{r} \subset \Hk_{\Cal{M}^{\all}}^{r}$ by the condition that $t_0 $ base changes to an injective map along every geometric point $\ol{s} \rightarrow S$ (equivalently, every $t_i$ has this property). Let $\pr_{i}: \Hk^{r}_{\cM}\to \cM$ (resp. $\pr_i^{\all} \co \Hk^r_{\cM^{\all}} \to \cM^{\all}$) be the map recording $(\cE, \cF_{i}, h_{i}, t_{i})$, for $0\le i\le r$.

\end{defn}


\begin{defn}[Shtukas for Hitchin stacks]
\source{higher-siegel-weil-unitary-nonsingular-terms}
\notready\label{defn: hitchin shtuka}
\source{higher-siegel-weil-unitary-nonsingular-terms}
For $r \geq 0$, we define $\Sht_{\Cal{M}^{\all}}^{r}$ as the fibered product 
\begin{equation}\label{eq: shtuka all as intersection}
\source{higher-siegel-weil-unitary-nonsingular-terms}
\begin{tikzcd}
\Sht_{\Cal{M}^{\all}}^{r}  \ar[r] \ar[d] & \Hk_{\Cal{M}^{\all}}^{r} \ar[d, "{(\pr_0^{\all}, \pr_r^{\all})}"] \\
\Cal{M}^{\all} \ar[r, "{(\Id, \Frob)}"] & \Cal{M}^{\all}  \times \Cal{M}^{\all}
\end{tikzcd}
\end{equation}
and the open substack $\Sht_{\Cal{M}}^{r} \subset \Sht_{\Cal{M}^{\all}}^{r}$ as the fibered product 
\begin{equation}\label{eq: shtuka as intersection}
\source{higher-siegel-weil-unitary-nonsingular-terms}
\begin{tikzcd}
\Sht_{\Cal{M}}^{r}  \ar[r] \ar[d] & \Hk_{\Cal{M}}^{r} \ar[d, "{(\pr_0, \pr_r)}"] \\
\Cal{M} \ar[r, "{(\Id, \Frob)}"] & \Cal{M}  \times \Cal{M}
\end{tikzcd}
\end{equation}
Explicitly, the stack $\Sht_{\Cal{M}^{\all}}^{r}$ parametrizes diagrams of the form below, with notation as in Definition \ref{defn: hitchin hecke},
\begin{equation}\label{eq: sht_M point}
\source{higher-siegel-weil-unitary-nonsingular-terms}
\begin{tikzcd}
\Cal{E} \ar[d, "t_0"] \ar[r, equals] &  \Cal{E} \ar[r, equals]  \ar[d, "t_1"]& \ldots \ar[r, equals] & \Cal{E} \ar[r, "\sim"] \ar[d, "t_{r}"] & \ft \Cal{E} \ar[d, "\ft t_0"] \\
\Cal{F}_0 \ar[r, dashed] & \Cal{F}_1 \ar[r, dashed]  &  \ldots \ar[r, dashed] &  \Cal{F}_{r} \ar[r, "\sim"] & \ft \Cal{F}_0  
\end{tikzcd}
\end{equation}
and $\Sht_{\Cal{M}}^{r}$ is the open substack where $t_0$ base changes to an injective map along every geometric point $\ol{s} \rightarrow S$ (equivalently, the same property holds for every $t_i$). 

%We let $[\Sht_{\Cal{M}_{H_1,H_2}}^{r}] \in \Ch_*(\Sht_{\Cal{M}_{H_1,H_2}}^{r})$ denote the class arising from \eqref{eq: shtuka as intersection} via the refined Gysin homorphism \cite[\S A.1.4]{YZ}. 
\end{defn}

%\mnote{Y: we have to define a virtual fundamental on  $\Hk^{r}_{\cM_{H_{1},H_{2}}}$. This is subtle since it may have junk components. This is circumvented by tracing it back to $\Hk^{1}_{\cM}$, which is smooth. See discussion in \S\ref{ss:ShtM cycle}.}


In particular, $\Cal{E} \xrightarrow{\sim} \ft \Cal{E}$ is a shtuka with no legs, exhibiting $\Cal{E}$ as arising from a rank $m$ vector bundle on $X'$, i.e. coming from $\Bun_{\GL'_m}(k)$. Note that in \eqref{eq: sht_M point}, $\cE$ is not fixed, so the automorphisms of $\cE$ are present in the functor of points of $\Sht_{\Cal{M}^{\all}}^{r}$. Therefore, if we define
\[
\ol{\Cal{Z}}_{\Cal{E}}^{r} :=  [\Cal{Z}_{\Cal{E}}^{r} /  (\Aut(\cE)(k)],
\]
then $\Sht_{\Cal{M}^{\all}}^{r}$ decomposes as a disjoint union of special cycles
\begin{equation}\label{eq: hitchin shtuka decomposition E}
\source{higher-siegel-weil-unitary-nonsingular-terms}
\Sht_{\Cal{M}^{\all}}^{r} = \coprod_{\Cal{E} \in \Bun_{\GL_{m}'}(k)} \ol{\Cal{Z}}_{\Cal{E}}^{r}.
\end{equation}
This decomposition can be refined.
%There are projections $\pr_i \co \Hk_{\Cal{M}}^{r}  \rightarrow \Cal{M}$ for $i=0, 1, \ldots, r$. 
The compositions $f \circ \pr_i \co \Hk_{\Cal{M}^{\all}}^{r} \rightarrow \Cal{A}^{\mrm{all}}$ all coincide, and they induce a map 
\begin{equation}\label{eq: hitchin shtuka decomposition A}
\source{higher-siegel-weil-unitary-nonsingular-terms}
\Sht_{\Cal{M}^{\all}}^{r} \rightarrow \Cal{A}^{\mrm{all}}(k).
\end{equation}
Here $\cA^{\mrm{all}}(k)$ is the groupoid of pairs $\{\cE, \cE \xrightarrow{a} \sigma^* \cE^\vee\}$. Let us write $\Aut(a) := \Aut(\cE \xrightarrow{a} \sigma^* \cE^\vee)\subset \Aut(\cE)$. Then each $\cE \xrightarrow{a} \sigma^* \cE^\vee$ defines a map $a: \Spec k \to B \Aut(a)(k)\incl \cA^{\all}(k)$, the latter map being an open-closed inclusion, and we have a Cartesian square
\begin{equation}\label{eq: hitchin shtuka decomposition E and A}
\source{higher-siegel-weil-unitary-nonsingular-terms}
\begin{tikzcd}
\cZ_{\cE}^r(a) \ar[r] \ar[d] & \Sht_{\Cal{M}^{\all}}^{r}  \ar[d, "f"] \\
\Spec k \ar[r, "a"] & \cA^{\mrm{all}}(k) 
\end{tikzcd}
\end{equation}
In particular, $\cZ_{\cE}^r(a)$ is a finite \'{e}tale cover of an open-closed substack of $\Sht_{\Cal{M}^{\all}}^{r}$ isomorphic to $[\cZ_{\cE}^r(a)/\Aut(a)(k)]$. 
Similarly, restricting to the injective locus (see Definition \ref{def:Z(0)}) we have a 
Cartesian square
\begin{equation}\label{eq: hitchin shtuka decomposition E and A inj}
\source{higher-siegel-weil-unitary-nonsingular-terms}
\begin{tikzcd}
\cZ_{\cE}^r(a)^{\circ} \ar[r] \ar[d] & \Sht_{\Cal{M}}^{r}  \ar[d, "f"] \\
\Spec k \ar[r, "a"] & \cA^{\mrm{all}}(k) 
\end{tikzcd}
\end{equation}
If $m=n$ and $a\in \cA(k)$, then we have $\cZ_{\cE}^r(a)^{\circ}=\cZ_{\cE}^r(a)$ and we can replace $ \cA^{\mrm{all}}(k)$ by $ \cA(k)$ in the above diagram (see Remark \ref{r:M to A}). 


%\[
%\Sht_{\Cal{M}^{\all}}^{r} = \coprod_{\Cal{E} \in \Bun_{m,X'}(k)}  \coprod_{a \in \cA^{\rm all}_{\Cal{E}}(k)} \ol{\Cal{Z}}_{\Cal{E}}^{r}(a)
%\]
%where $\cA^{\rm all}_{\cE}$ is the fiber of $\cA^{\rm all} \rightarrow \Bun_{m,X'}$ over $\cE \in \Bun_{m, X'}(k)$. (Clearly the $k$-points of $\cA^{\rm all}_{\cE}$ coincide with $\cA^{\rm all}_{\cE}(k)$ as defined in \S \ref{ssec: sht index by a}.) This induces a decomposition of the open substacks,
%\begin{equation}\label{eq: hitchin shtuka decomposition E and A}
%\Sht_{\Cal{M}}^{r} = \coprod_{\Cal{E} \in \Bun_{m,X'}(k)}  \coprod_{a \in \cA^{\rm all}_{\Cal{E}}(k)} \ol{\Cal{Z}}_{\Cal{E}}^{r}(a)^\circ.
%\end{equation}







\subsection{From vector bundles to torsion sheaves}\label{ss:vb to tor}
\source{higher-siegel-weil-unitary-nonsingular-terms}
For the rest of the section, we concentrate on the case $m=n$. In this case, we will relate $\cM = \cM(n,n)$ to the moduli stack of Hermitian torsion sheaves introduced in \S\ref{ss:Herm}.  We introduce the following abbreviated notations.

\begin{defn}
\source{higher-siegel-weil-unitary-nonsingular-terms}
\notready Let $d\in \Z_{\ge0}$.
\begin{enumerate}
\item Let $\cM_{d} = \cM(n,n)_d$ be the open-closed substack of $\cM = \cM(n,n)$ consisting of $(\cE\xr{t}\cF)$ where $d=\frac{\deg\cE^\vee-\deg(\cE)}{2}=-\chi(X',\cE)$.  
\item Let $\cA_{d}  = \cA(n)_d$ be the open-closed substack of $\cA = \cA(n)$ consisting of $(\cE,a)$ where $d=\frac{\deg\cE^\vee-\deg(\cE)}{2}=-\chi(X',\cE)$.
\end{enumerate}
\end{defn}

By Remark \ref{r:M to A}, the Hitchin map for $\cM ^{\rm all}= \cM^{\rm all}(n,n)$ restricts to a map 
\begin{equation}
f_{d}: \cM_{d}\to \cA_{d}.
\end{equation}
When $d$ is understood, we abbreviate $f$ for $f_{d}$.


Recall that  $\Herm_{2d}=\Herm_{2d}(X'/X)$ is the moduli stack of length $2d$ torsion coherent sheaves $\cQ$ on $X'$ equipped with a Hermitian structure $h_{\cQ} \co \cQ \xrightarrow{\sim} \s^{*}\cQ^{\vee}$, where $\cQ^{\vee} := \ul{\Ext}^1(\cQ, \omega_{X'})$ such that $\sigma^* h_{\cQ}^{\vee}=h_{\cQ}$.  Alternatively we may think of $h_{\cQ}$ as the datum of a perfect pairing 
\begin{equation}
h'_{\cQ}: \cQ\ot_{\cO_{X'}}\s^{*}\cQ\to \omega_{X',F'}/\omega_{X'}
\end{equation}
where $\omega_{X',F'}$ is the constant Zariski sheaf of rational differential form on $X'$. The Hermitian condition is equivalent to $h'_{\cQ}(u,v)=\s^{*}h'_{\cQ}(v,u)$ for local sections $u,v$ of $\cQ$.
 
In \S\ref{ss:HSpr stalk} we have also introduced the moduli stack $\Lagr_{2d}=\Lagr_{2d}(X'/X)$ classifying $(\cQ,h_{\cQ}, \cL)$ where $(\cQ,h_{\cQ})\in \Herm_{2d}$ and $\cL\subset \cQ$ is a Lagrangian subsheaf. 


There is a canonical map $g: \Cal{A}_d \rightarrow \Herm_{2d}$ sending $(\Cal{E},a)$ to the torsion sheaf $\cQ=\sigma^* \Cal{E}^{\vee}/\Cal{E}$ together with the Hermitian structure $h_{\cQ}$ defined in \S \ref{ssec: density function for torsion sheaves}. 
%\mnote{Y: this part may be deleted.} 
%To check $h_{\cQ}$ is Hermitian we pass to its equivalent datum $h'_{\cQ}$. The Hermitian pairing $h: \cE\ot \s^{*}\cE^{\vee}\to \omega_{X'}$ induces a rational pairing
%\begin{equation}
%h'_{F'}: \s^{*}\cE^{\vee}\ot \cE^{\vee}\to \omega_{X',F'}
%\end{equation}
%and passes to the quotient
%\begin{equation}
%h'_{\cQ}: \cQ\ot \s^{*}\cQ\to \omega_{X',F'}/\omega_{X'}
%\end{equation}
%which is the equivalent datum of $h_{\cQ}$. Since $h'$ and $h'_{F'}$ are Hermitian, so is $h'_{\cQ}$. 



%Let $\Lagr_d = \{(Q,h_Q, \Cal{L} \text{ Lagrangian}\subset Q\}$. 


We have a map $g_{\cM}: \Cal{M}_d \rightarrow \Lagr_{2d}$ sending $(\Cal{E} \xr{t} \Cal{F}) \in \Cal{M}_d$ to the torsion Hermitian sheaf $(\cQ=\sigma^* \Cal{E}^{\vee}/\Cal{E}, h_{\cQ})$ constructed above together with the Lagrangian  $\cL=\Cal{F}/\cE$. 


\begin{lemma}
\source{higher-siegel-weil-unitary-nonsingular-terms}
\notready\label{lem: hitchin cartesian} The maps defined above fit into a Cartesian diagram 
\source{higher-siegel-weil-unitary-nonsingular-terms}
\begin{equation}\label{eq: hitchin cartesian}
\source{higher-siegel-weil-unitary-nonsingular-terms}
\begin{tikzcd}
\Cal{M}_d \ar[d, "f"] \ar[r,"g_{\cM}"]  & \Lagr_{2d} \ar[d, "\upsilon_{2d}"] \\
\Cal{A}_d \ar[r, "g"] & \Herm_{2d}
\end{tikzcd}
\end{equation}
\end{lemma}

\begin{proof} Given $a: \Cal{E} \rightarrow \sigma^* \Cal{E}^{\vee}$ that is injective, the datum of a subsheaf $\cL\subset \sigma^* \Cal{E}^{\vee}/\Cal{E}$ is the same as a coherent sheaf $\cF$ such that $\Cal{E} \subset \Cal{F} \subset \sigma^* \Cal{E}^{\vee}$. It is easy to see that $\cL$ is Lagrangian if and only if $\cF$ is self-dual under the Hermitian map $a$.
\end{proof}



\begin{cor}
\source{higher-siegel-weil-unitary-nonsingular-terms}
\notready\label{c: hitchin proper}
\source{higher-siegel-weil-unitary-nonsingular-terms}
The Hitchin fibration $f_{d} \co \Cal{M}_{d} \rightarrow \Cal{A}_{d}$ is proper.
\end{cor}
\begin{proof}Apply Lemma \ref{lem: hitchin cartesian} and the fact that $\upsilon_{2d}$ is proper.
\end{proof}



\begin{lemma}
\source{higher-siegel-weil-unitary-nonsingular-terms}
\notready\label{l:Lag small}
\source{higher-siegel-weil-unitary-nonsingular-terms}
The map $\upsilon_{2d} \co \Lagr_{2d} \rightarrow \Herm_{2d}$ is small.
\end{lemma}
\begin{proof}
The map $\pi^{\Herm}_{2d}$ from \S\ref{ss:Herm Spr} factors as 
\begin{equation}\label{factor pi Herm}
\source{higher-siegel-weil-unitary-nonsingular-terms}
\pi^{\Herm}_{2d}: \wt\Herm_{2d}\xr{\l_{2d}}\Lagr_{2d}\xr{\upsilon_{2d}}\Herm_{2d}.
\end{equation}
Since $\l_{2d}$ is surjective and $\pi^{\Herm}_{2d}$ is small by Proposition \ref{p:SH Spr}, we get the desired statement.
\end{proof}

%\begin{proof}
%By Lemma \ref{l:Herm local}, $\Herm_{2d}$ is \'etale locally isomorphic to $[\so_{2d}/\Og_{2d}]$. Under such an \'etale chart, the map $\upsilon_{2d}: \Lagr_{2d}\to \Herm_{2d}$ is locally isomorphic to the map $\upsilon_{2d}: [\frp/P]\to [\so_{2d}/\Og_{2d}]$ where $P\subset \Og_{2d}$ is the stabilizer of a Lagrangian subspace (there are two families of Lagrangians under $\SO_{2d}$-action but only one under $\Og_{2d}$), and $\frp=\Lie P$. It is well-known that $\upsilon_{2d}$ is a small map \footnote{Reference: Lusztig or Borho-MacPherson}. Therefore $\upsilon_{2d}$ is also small. 
%\end{proof}

\subsubsection{Proof of Proposition \ref{prop: KR cycle proper}}\label{sssec: complete proof}
\source{higher-siegel-weil-unitary-nonsingular-terms}
Let $\Lagr(\cQ)$ be the moduli space of Lagrangian subsheaves of $\cQ$. Let $\Hk^{r}_{\Lagr(\cQ)}$ be its Hecke version, classifying points $\{x'_{i}\}_{1\le i\le r}$ of $X'$ and chains of Lagrangian subsheaves of $\cQ$
\begin{equation}
\begin{tikzcd}
\cL_{0} \ar[r, dashed, "{f_{1}}"]   & \cL_{1} \ar[r,dashed, "{f_{2}}"]  & \cdots \ar[r, dashed, "{f_{r}}"]  & \cL_{r}
\end{tikzcd}
\end{equation}
where the dashed arrow $f_{i}$ are modifications at $x'_{i}\cup \s(x'_{i})$, similar to those in Definition \ref{defn: Hk U(n)}. There is a natural map $\cZ^{r}_{\cE}(a)\to \Hk^{r}_{\Lagr(\cQ)}$ sending a point $(\{x'_{i}\}, \{t_{i}: \cE\to \cF_{i}\})$ of $\cZ^{r}_{\cE}(a)$ to the collection of (necessarily Lagrangian) subsheaves $\ov \cF_{i}=\coker(t_{i})\subset\cQ=\s^{*}\cE^{\vee}/\cE$. This map fits into a Cartesian diagram
\begin{equation}\label{ZLag}
\source{higher-siegel-weil-unitary-nonsingular-terms}
\xymatrix{\cZ^{r}_{\cE}(a)\ar[r]\ar[d] & \Hk^{r}_{\Lagr(\cQ)} \ar[d] \\
\Lagr(\cQ)\ar[r]^-{(\Id,\Frob)} & \Lagr(\cQ)\times \Lagr(\cQ)}
\end{equation}
Now both $\Lagr(\cQ)$ and $\Hk^{r}_{\Lagr(\cQ)}$ are proper schemes over $k$, hence the same is true for  $\cZ^{r}_{\cE}(a)$.  The diagram also makes it clear that $\cZ^{r}_{\cE}(a)$ only depends on $(\cQ,\ov a)$.


\subsection{Smoothness}


\begin{lemma}
\source{higher-siegel-weil-unitary-nonsingular-terms}
\notready\label{l:Md smooth}  The map 
\source{higher-siegel-weil-unitary-nonsingular-terms}
\begin{equation}
\b_{d}: \cM_{d}\to \Coh_{d}(X')\times \Bun_{U(n)}
\end{equation}
sending $(\cE\xr{t}\cF,h)$ to $(\coker(t), (\cF,h))$ is smooth of relative dimension $dn$. In particular, $\cM_{d}$ is smooth of pure dimension $dn+n^{2}(g-1)$. 
\end{lemma} 
\begin{proof}
Consider the stack $\cM'_{d}$ classifying $(\cT,\cF,h, s)$ where $\cT\in \Coh_{d}(X')$, $(\cF,h)\in \Bun_{U(n)}$ and an $\cO_{X'}$-linear map $s:\cF\to \cT$. Let $\b'_{d}: \cM'_{d}\to \Coh_d(X') \times \Bun_{U(n)}$ be the natural map. Due to the vanishing of $\Ext^{1}(\cF,\cT)$, $\b'_{d}$ exhibits $\cM'_{d}$ as a vector bundle over $\Coh_d(X') \times \Bun_{U(n)}$ of rank $dn=\dim\Hom_{X'}(\cF,\cT)$. Now $\cM_{d}$ is the open substack of $\cM'_{d}$ where $s$ is surjective. Therefore $\b_{d}$ is also smooth of relative dimension $dn$. 
\end{proof}



\begin{prop}
\source{higher-siegel-weil-unitary-nonsingular-terms}
\notready\label{prop:Ad Herm smooth}
\source{higher-siegel-weil-unitary-nonsingular-terms}
The map $g \co \Cal{A}_d \rightarrow \Herm_{2d}$ is smooth. 
\end{prop}
%Using Lemma \ref{lem: hitchin cartesian}, Theorem \ref{th:small} follows from combining the two results below. 
%\begin{prop}\label{prop: 1}
%The map $\upsilon_{2d} \co \Lagr_{2d} \rightarrow \Herm_{2d}$ is small.
%\end{prop}
%\begin{proof}[Proof of Proposition \ref{prop: 1}] 
%
%\end{proof}


%We claim that \'{e}tale locally, $\Herm_d$ is isomorphic to $[\mf{gl}_d/_{\Ad} \GL_d]$. The point is that \'{e}tale locally the double cover is split, in which case $\Herm_d \cong \mrm{Coh}_d$ by the map 
%$\Cal{T} \sqcup \Cal{T}^{\vee} \leftarrow \Cal{T}$. \mnote{more justification, probably much earlier}
%
%The local structure of $\mrm{Coh}_d$ is well-known: \'{e}tale locally $X \cong \A^1$, and $\mrm{Coh}_d(\A^1) \cong \mf{gl}_d/_{\Ad} \GL_d$ since a length $d$ coherent module over $k[t]$ is just a vector space with an endomorphism; the map sends the endomorphism to its matrix. 
%
%Similarly, \'{e}tale locally $\Lagr_d$ is isomorphic to $\coprod_{i=0}^d \Coh_{d,i}(X)$ where $\Coh_{d,i}(X)$ parametrizes a length $i$ torsion subsheaf $\Cal{T}' \subset \Cal{T}$. This is because in the split case $X' = X \sqcup X$, $Q$ must be of the form $\Cal{T} \oplus \Cal{T}^{\vee}$ and $L$ must be of the form $\Cal{T}' \oplus (\Cal{T}')^{\perp}$. 
%
%For $X=\A^1$, the map $\wt{\mf{gl}}_d^i \rightarrow \mf{gl}_d$ is a partial Springer resolution, with 
%\[
%\wt{\mf{gl}}_d^i = \{(A \in \mf{gl}_d, V \subset \A^d) \co A V \subset V\}.
%\]

\begin{proof}
We have a map $s^{\Lagr}_{2d} \co \Lagr_{2d}\to X'_{d}$ sending $(\cQ,h_{\cQ},\cL)$ to the divisor of $\cL$. Recall also the map $s^{\Herm}_{2d}: \Herm_{2d}\to X_{d}$ sending $(\cQ,h_{\cQ})$ to the descent of the divisor of $\cQ$ to $X$.

Recall in \S\ref{ss:HSpr stalk} we introduced the open subset 
\[
(X_d')^{\dm} = \{D' \subset X' \co D' \cap \sigma(D') = \vn\} \subset X_d'.
\]
Let $\Lagr_{2d}^{\dm} \subset \Lagr_{2d}$ and $\Cal{M}_d^{\dm} \subset \Cal{M}_d$ be the preimages of $(X'_{d})^{\dm}$ under the maps $s^{\Lagr}_{2d}$ and $s^{\Lagr}_{2d}\c g_{\cM}$. 

%We have a commutative diagram 
%\[
%\begin{tikzcd}
%\Cal{M}_d \ar[d, "f_d"]  \ar[r] & \Lagr_{2d} \ar[d] \ar[r] & (X')_d \ar[d] \\
%\Cal{A}_d \ar[r] & \Herm_{2d} \ar[r] & X_{d}
%\end{tikzcd}
%\]

We claim that both squares in the diagram 
\[
\begin{tikzcd}
\Cal{M}_d^{\dm} \ar[r] \ar[d, "f_d"] & \Lagr_{2d}^{\dm} \ar[d] \ar[r] & (X'_d)^{\dm}\ar[d] \\
\Cal{A}_d \ar[r,"g"] & \Herm_{2d} \ar[r] & X_d
\end{tikzcd}
\]
are Cartesian. The left square is Cartesian by definition. Now we show that the right square is Cartesian. Let $(\cQ ,h_{\cQ})\in \Herm_{2d}$,  $D' \in (X_d')^{\dm}$ lying over $D=s^{\Herm}_{2d}(\cQ,h_{\cQ})$. Since  $D'\cap\sigma(D')=\vn$, there is a unique Lagrangian subsheaf $\cL \subset \cQ$ supported on the support of $D'$, namely $\cL=\cQ|_{\supp D'}$. This gives the unique point $(\cQ,h_{\cQ}, \cL)\in \Lagr^{\dm}_{2d}$ mapping to $(\cQ,h_{\cQ})\in \Herm_{2d}$ and $D'\in (X'_{d})^{\dm}$. 

Note that the map $(X'_d)^{\dm} \rightarrow X_d$ is faithfully flat: it is clearly surjective, and the map $\nu_{d}: X'_d \rightarrow X_d$ is a finite morphism between smooth schemes, hence flat. We will show that $\Cal{M}_d^{\dm} \rightarrow \Lagr_d^{\dm}$ is smooth. By fppf descent it then follows that $\Cal{A}_d \rightarrow \Herm_d$ is also smooth. 

Recall from \S\ref{ss:HSpr stalk} that $\e'_{d}: \Lagr_{2d}\to \Coh_{d}(X')$ (recording only $\cL$) restricts to an isomorphism $\Lagr_{2d}^{\dm}\isom \Coh_{d}(X')^{\dm}:=\Coh_d(X')|_{(X_d')^{\dm}}$. Therefore it suffices to show that the composition $\cM_{d}^{\dm}\xr{g_{\cM}}\Lagr_{2d}^{\dm}\xr{\e'_{d}} \Coh_{d}(X')^{\dm}$ is smooth. This follows from the smoothness of $\cM_{d}\to \Coh_{d}(X')$ proved in Lemma \ref{l:Md smooth}.
\end{proof}

%Note that sending $(L \subset Q) \in \Lagr_d^{\dm}$ necessarily has $Q = L \oplus \sigma^* (L^{\vee})$. Therefore the map $(L \subset Q) \mapsto L$ defines an isomorphism  $ \Lagr_d^{\dm}\cong \Coh_d(X')|_{(X_d')^{\dm}}$. 
%
%The composite $\Cal{M}_d^{\dm} \rightarrow \Lagr_d^{\dm} \cong \Coh_d(X')|_{(X_d')^{\dm}}$ is the localization over $(X_d')^{\dm}$ of the map $\Cal{M}_d \rightarrow \Coh_d(X')$ sending $(\Cal{E} \rightarrow \Cal{F}) \in \Cal{M}_d$ to $\coker(\Cal{E} \rightarrow \Cal{F}) \in \Coh_d(X')$. So it suffices to show that the latter is smooth. 

%Recall that $\Cal{M}_d^{\dm}$ parametrizes maps of vector bundles $\Cal{E} \rightarrow (\Cal{F},h)$ on $X'$ of colength $d$, and where $\Cal{F}$ has also a Hermitian structure. 


%Let $\cM'_{d}$ be the component of $\cM_{\GL(n)',\GL(n)'}$ classifying injective maps $(\cE\to \cF)$ of rank $n$ vector bundles over $X'$ with $\deg\cE=-d$ and $\deg\cF=0$.  Tautologically, we have a cartesian diagram 
%\[
%\begin{tikzcd}
%\Cal{M}_d \ar[r] \ar[d] & \cM'_{d}  \ar[d] & (\Cal{E} \rightarrow \Cal{F}) \ar[d] \\
%\Bun_{U(n)} \ar[r] & \Bun_{n,X'} &  \Cal{F} 
%\end{tikzcd}
%\]
%We claim that the map $h: \cM'_{d} \rightarrow \Coh_d(X') \times \Bun_{n,X'}$ sending $(\cE\to\cF)$ to $(\cF/\cE, \cE)$ is smooth of relative dimension $dn$. \footnote{It is not surjective; the image is $\Coh_d^{(n)} (X') \times \Bun_n(X')$ where $\Coh_d^{(n)} (X')  \subset \Coh_d(X')$ is the open substack parametrizing coherent sheaves admitting a surjection from $\Cal{O}_{X'}^{\oplus n}$.} By the above Cartesian diagram, this implies our claim that $\cM_{d}\rightarrow \Coh_d(X')$ is smooth. 
%
%To check the smoothness of $h$, consider the stack $\cH\to \Coh_d(X') \times \Bun_{n,X'}$ classifying $(\cL,\cF, s)$ where $\cL\in \Coh_{d}(X')$, $\cF\in \Bun_{n,X'}$ and an $\cO_{X'}$-linear map $s:\cF\to \cL$. Due to the vanishing of $\Ext^{1}(\cF,\cL)$, $\cH$ is a vector bundle over   $\Coh_d(X') \times \Bun_{n,X'}$ of rank $dn$. The stack $\cM'_{d}$ is the open substack of $\cH$ where $s$ is surjective. Therefore $h$ is smooth of relative dimension $dn$. This finishes the proof.


%Now the fiber of $h$ over $(\cL, \Cal{F})\in \Coh_d(X') \times \Bun_{n,X'}$ is the open (possibly empty) subscheme of $\Hom(\cF,\cL)$ consisting of surjective maps. \mrm{Surj}(\Cal{F}, L)$ which is an open () substack of $\Hom(\Cal{F}, L)$. We need to check flatness. 

%We analyze by strata. The ``generic'' strata is when $Q$ is of the form $k_{x_1' \oplus \ldots \oplus k_{x_d'}} \oplus k_{x_1''} \oplus \ldots \oplus k_{x_d''}$, where by convention $x'' = \sigma(x')$. 

%In this case, there are $2^d$ choices for the self-dual $\Cal{E} \hookrightarrow \Cal{F} \hookrightarrow \sigma^* \Cal{E}^{\vee}$. Indeed, $\Cal{M}_d \rightarrow \Cal{A}_d$ will be generically a $2^d$-cover. Consider the diagram 
%\[
%\begin{tikzcd}
%& \{ \Cal{E} \hookrightarrow \Cal{F} \hookrightarrow \sigma^* \Cal{E}^{\vee} \} \ar[dl, "\text{finite}"] \ar[dr] \\
%g^{-1}(S_Q) = \{\Cal{E} \rightarrow \sigma^* \Cal{E}^{\vee}\} & & \Bun_{U(n)}
%\end{tikzcd}
%\]
%So $\dim g^{-1}(S_{Q}) = \dim \Bun_{U(n)} + (n-1) d$. 

%Next consider the stratum where $Q$ is concentrated on a pair of conjugate points $x', x''$. Assume $n \mid d$ and suppose the Jordan blocks are $(\frac{d}{n}, \ldots, \frac{d}{n})$. (This corresponds to $\Cal{O}/t^{d/n} \oplus \ldots \oplus \Cal{O}/t^{d/n}$.) In this case, $h^{-1}(S_Q)$ consists of maps $h \co \Cal{E} \rightarrow \sigma^* \Cal{E}^{\vee}$ such that 
%\[
%\Cal{E} \hookrightarrow \Cal{E}(\frac{d}{n}x') \xrightarrow{\sim} \sigma^* \Cal{E}^{\vee}(-\frac{d}{n} x'') \hookrightarrow \sigma^* \Cal{E}^{\vee}.
%\]
%That is, to a point in $h^{-1}(S_Q)$ we can associated a $U(n)$-bundle and vice versa, so $h^{-1}(S_Q) \xrightarrow{\sim} \Bun_{U(n)}$ via $\Cal{E} \mapsto \Cal{E}(\frac{d}{n}x')$. 

%To check that the dimensions shake out correctly, let $P$ be the parabolic associated to the partition $(\frac{d}{n}, \ldots, \frac{d}{n})$. The generic  element of $n_P$ has that type, and letting $\Cal{O}$ denote its orbit closure, a variant of the springer resolution gives $T^*G/P \rightarrow \Cal{O}$. So 
%\[
%\dim \Cal{N} - \dim \Cal{O} = \dim T^* (G/B)  - \dim T^*(G/P) = 2 \dim (P/B).
%\]
%Since $P/B$ is the flag variety for $\GL_n^{d/n}$, this is $\frac{d}{n} 2 \frac{n(n-1)}{2} = d(n-1)$, which is what we were looking for. 


\begin{cor}
\source{higher-siegel-weil-unitary-nonsingular-terms}
\notready\label{c:small} 
\source{higher-siegel-weil-unitary-nonsingular-terms}
The Hitchin fibration $f \co \cM_d \rightarrow \cA_d$ is small. The complex $Rf_{*}\Qlbar$ is a shifted perverse sheaf that is the middle extension from any dense open substack of $\cA_{d}$.
\end{cor}
\begin{proof}
By the smoothness of $g$ in Proposition \ref{prop:Ad Herm smooth} and the Cartesian diagram in Lemma \ref{lem: hitchin cartesian}, the smallness of $f$ follows from that of $\upsilon_{2d}: \Lagr_{2d}\to \Herm_{2d}$, which is proved in Lemma \ref{l:Lag small}.
\end{proof}



%\begin{prop}\label{p:dim Hk image} Dimension of $\Hk^{r}_{\cM}$...
%\end{prop}
%\begin{proof}
%For a geometric point $(\{x_{i}\}, \{\cE\xr{t_{i}}\cF_{i}\})\in  \Hk^{r}_{\cM^{d}_{\GL(n)', U(n)}}$, the obstruction to infinitesimally deforming within the fiber over $(\{x_{i}\}, \{\cF_{i}\})\in \Hk^{r}_{U(n)}$ lies in $\upH^{2}(X', \cK)$ where $\cK$ is the following complex placed in degrees $0$ and $1$
%\begin{equation}\label{deformation HkM}
%\un\End(\cE)\xr{\th} \un\Hom(\cE, \cF_{\bu}).
%\end{equation}
%Here, using notation from Remark \ref{rem: formulation of upper/lower}, $\un\Hom(\cE, \cF_{\bu})$ is the kernel of 
%\begin{equation}
%\d: \bigoplus_{i=0}^{r}\un\Hom(\cE, \cF_{i})\to \bigoplus_{i=1}^{n}\un\Hom(\cE, \cF_{i-1/2}^{\sh})
%\end{equation}
%given by $(\ph_{0},\cdots,\ph_{r})\mapsto (\ph_{0}-\ph_{1},\ph_{1}-\ph_{2},\cdots, \ph_{r-1}-\ph_{r})$. The map $\th$ in \eqref{deformation HkM}  is given by $\a\mapsto (t_{0}\c\a,\cdots, t_{r}\c\a)$. Restricting \eqref{deformation HkM} to the generic point $\y$ of $X'$,  it becomes $\End(\cE_{\y})\to \Hom(\cE_{\y}, \cF_{0,\y})$ given by $\a\mapsto t_{0,\y}\c\a$. Since $t_{0,\y}$ is an isomorphism by assumption, $\th_{\y}$ is an isomorphism, hence $\coker(\th)$ is torsion. Therefore $\upH^{2}(X',\cK)\cong \upH^{1}(X',\coker(\th))=0$. This proves the smoothness of $\Hk^{r}_{\cM^{d}_{\GL(n)', U(n)}}$ over $\Hk^{r}_{U(n)}$.  The relative dimension is $-\chi(X', \cK)$. Now $\un\End(\cE)$ has degree $0$ and rank $n^{2}$.  
%\end{proof}




\subsection{Cycle class from Hitchin shtukas}\label{ss:ShtM cycle}
\source{higher-siegel-weil-unitary-nonsingular-terms}

In this subsection we take $m=n$, so $\cM = \cM(n,n)$ and $\cA = \cA(n)$. Now consider the Hitchin shtukas for $\cM_{d}\subset \cM$. Let $N=\dim \cM_{d}$. By Corollary \ref{c:small}, $\dim \cA_{d}=N$. The Cartesian diagram \eqref{eq: shtuka as intersection} restricts to a Cartesian diagram
\begin{equation}\label{def ShtMd}
\source{higher-siegel-weil-unitary-nonsingular-terms}
\xymatrix{\Sht_{\cM_{d}}^{r}\ar[d] \ar[r] & \Hk^{r}_{\cM_{d}}\ar[d]^{(\pr_{0},\pr_{r})}\\
\cM_{d}\ar[r]^-{(\Id, \Frob)} & \cM_{d}\times\cM_{d}}
\end{equation}
We would like to define a $0$-cycle class on $\Sht_{\cM_{d}}^{r}$ as the Gysin pullback of a cycle on $\Hk^{r}_{\cM_{d}}$ along the Frobenius graph of $\cM_{d}$. Although the virtual dimension of $\Hk^{r}_{\cM_{d}}$ is the same as $\dim\cM_{d}$, its actual dimension may be larger. For this reason we have to define in a roundabout way a virtual fundamental cycle on $\Hk^{r}_{\cM_{d}}$ of the virtual dimension by relating it to $\Hk^{1}_{\cM_{d}}$, which we show is smooth below. 

\begin{lemma}
\source{higher-siegel-weil-unitary-nonsingular-terms}
\notready\label{l:HkMd sm} The stack $\Hk^{1}_{\cM_{d}}$ is smooth and equidimensional of the same dimension as $\cM_{d}$.
\source{higher-siegel-weil-unitary-nonsingular-terms}
\end{lemma}
\begin{proof}
Let $(x',\Cal{F}_0 \dashrightarrow \Cal{F}_1) \in \Hk_{U(n)}^1$. Let $\cF^{\flat}:=\cF^{\flat}_{1/2}=\Cal{F}_0 \cap \Cal{F}_1$ as in Definition \ref{defn: Hk U(n)}. 
%By means of the given rational isomorphism between $\Cal{F}_0$ and $\Cal{F}_1$, we may form the intersection $\cF^{\flat}=\Cal{F}_0 \cap \Cal{F}_1$. 
The generically compatible Hermitian structures on $\Cal{F}_0$ and $\Cal{F}_1$ equip this intersection with a Hermitian structure $h^{\flat}: \cF^{\flat} \inj \sigma^* (\Cal{F}^{\flat})^{\vee}$ whose cokernel has length $1$ at $x'$ and at $\sigma (x')$. We call such a bundle {\em almost Hermitian} with defect at the ordered pair of conjugate points $(x',\s(x'))$. Conversely, given $(\Cal{F}^{\flat},h^{\flat})$ almost Hermitian with defect at $(x', \s(x'))$, one can uniquely recover $\cF_{0}$ (resp. $\cF_{1}$) as the upper modification of $\cF^{\flat}$ at $x'$ (resp. $\s(x')$) inside $\s^{*}(\cF^{\flat})^{\vee}$.  

Let $\Bun^{\flat}_{U(n)}$ be the moduli stack parametrizing  $(x'\in X', \cF^{\flat}, h^{\flat})$ where $(\cF^{\flat}, h^{\flat})$ is an almost Hermitian bundle with defect at $(x',\sigma(x'))$. The discussion in the previous paragraph shows that there is an isomorphism $\Hk^{1}_{U(n)}\isom \Bun^{\flat}_{U(n)}$ over $X'$.  
Let $\cM^{\flat}_{d}$ be the moduli stack of $(x',\cE\xr{t}\cF^{\flat}, h^{\flat})$ where $(\cF^{\flat}, h^{\flat})$ is almost Hermitian with defect at $(x',\s(x'))$, $\cE$ is a vector bundle on $X'$ of rank $n$ and $\chi(X',\cE)=-d$, and $t$ is injective. Then we have an isomorphism $\Hk^{1}_{\cM_{d}}\cong \cM^{\flat}_{d}$.  


We have a natural map 
\begin{equation}
\b^{\flat}_{d}: \cM^{\flat}_{d}\to \Coh_{d-1}(X')\times \Bun^{\flat}_{U(n)}
\end{equation}
sending $(x',\cE\xr{t}\cF^{\flat}, h^{\flat})$ to $(\coker(t), (x',\cF^{\flat}, h^{\flat}))$. The same argument as Lemma \ref{l:Md smooth} shows that $\b^{\flat}_{d}$ exhibits $\cM^{\flat}_{d}$ as an open substack in a vector bundle of rank $n(d-1)$ over $\Coh_{d-1}(X')\times \Bun^{\flat}_{U(n)}$. Now $\dim \Coh_{d-1}(X')=0$ and $\dim \Bun^{\flat}_{U(n)}=\dim \Hk^{1}_{U(n)}=\dim \Bun_{U(n)}+n$ by Lemma \ref{lem: shtuka geometry}(1). Therefore $\Hk^{1}_{\cM_{d}}\cong \cM^{\flat}_{d}$ is of pure dimension $\dim \Bun_{U(n)}+dn$, which is the same as $\dim \cM_{d}$ by Lemma \ref{l:Md smooth}.
\end{proof}

%At this point we could define a virtual fundamental cycle on $\Hk^{r}_{\cM_{d}}$. Instead we shall directly define a $0$-cycle on $\Sht^{r}_{\cM_{d}}$ as follows. \mnote{the reader may find these two lines to be  non sequitur; could just delete them}

\begin{defn}
\source{higher-siegel-weil-unitary-nonsingular-terms}
\notready\label{def:Phi} For any stack $S$ over $k$ we define a morphism
\source{higher-siegel-weil-unitary-nonsingular-terms}
\begin{equation}
\Phi^{r}_{S}: S^{r+1}\to S^{2r+2}
\end{equation}
by the formula $\Phi^{r}_{S}(\xi_{0},\cdots, \xi_{r})=(\xi_{0},\xi_{1},\xi_{1},\xi_{2},\xi_{2},\cdots, \xi_{r-1}, \xi_r, \xi_r, \Frob(\xi_{0}))$. When $r$ is fixed in the context, we simply write $\Phi_{S}$.
\end{defn}

We rewrite $\Sht^{r}_{\cM_{d}}$ as the fiber product
\begin{equation}\label{rewrite ShtM}
\source{higher-siegel-weil-unitary-nonsingular-terms}
\xymatrix{\Sht^{r}_{\cM_{d}} \ar[r]\ar[d] & (\Hk^{1}_{\cM_{d}})^{r}\ar[d]^{(\pr_{0},\pr_{1})^{r} \times \Delta} \times \cM_d\\
(\cM_{d})^{r+1}\ar[r]^-{\Phi^{r}_{\cM_{d}}} & (\cM_{d})^{2r+2}=(\cM_{d})^{2r}\times(\cM_{d})^{2}}
\end{equation}
Here the vertical map $(\pr_{0},\pr_{1})^{r}$ sends $(h_{1},\cdots, h_{r})\in (\Hk^{1}_{\cM_{d}})^{r} $ to 
$(\pr_{0}(h_{1}), \pr_{1}(h_{1}),\cdots, \pr_{0}(h_{r}),\pr_{1}(h_{r}))\in (\cM_{d})^{2r}$, while $\Delta$ is the diagonal map. 

\begin{defn}
\source{higher-siegel-weil-unitary-nonsingular-terms}
\notready\label{def: hitchin shtukas class} We define a $0$-cycle classes $[\Sht^{r}_{\cM_{d}}]\in \Ch_{0}(\Sht^{r}_{\cM_{d}})$ as the image of the fundamental class of $(\Hk^{1}_{\cM_{d}})^{r} \times \cM_d$  (which is smooth of the same dimension as $(\cM_{d})^{r+1}$ by Lemma \ref{l:HkMd sm}) under the refined  Gysin map along $\Phi_{\cM_{d}}:(\cM_{d})^{r+1}\to (\cM_{d})^{2r+2}$ (which is defined since $\cM_{d}$ is smooth and equidimensional by Lemma \ref{l:Md smooth}; see \cite[Theorem 2.1.12(xi)]{Kr99}) 
\source{higher-siegel-weil-unitary-nonsingular-terms}
\begin{equation}
[\Sht^{r}_{\cM_{d}}]:=(\Phi^{r}_{\cM_{d}})^{!}[(\Hk^{1}_{\cM_{d}})^{r} \times \cM_d]\in \Ch_{0}(\Sht^{r}_{\cM_{d}}).
\end{equation}
\end{defn}

%Consider the map $\cA_{d}\to \Herm_{2d}\to X_{d}$ that maps $(\cE,a)$ to $D_{a}$. Let $\cA^{\c}_{d}$ be preimage of the open subset $X^{\c}_{d}$ of multiplicity-free divisors. Let $\cM^{\c}_{d}$ and $\Hk^{r,\c}_{\cM_{d}}$ be the preimages of $\cA^{\c}_{d}$ in $\cM_{d}$ and $\Hk^{r}_{\cM_{d}}$.
%
%The following lemma is easy to see.
%\begin{lemma}
%The map $\Hk^{r}_{\cM_{d}}\to \cA_{d}$ is  finite \'etale when restricted over $\cA^{\c}_{d}$.  
%\end{lemma}
%From this lemma, the fundamental class of the closure $\ov\Hk^{r,\c}_{\cM_{d}}$ in $\Hk^{r}_{\cM_{d}}$ defines a cycle class which we denote
%\begin{equation}
%[\Hk^{r, \c}_{\cM_{d}}]\in \Ch_{N}(\Hk^{r, \c}_{\cM_{d}}).
%\end{equation}
%
%\begin{defn} We define a $0$-cycle classes $[\Sht^{r}_{\cM_{d}}]\in \Ch_{0}(\Sht^{r}_{\cM_{d}})$ as the the image of the refined  Gysin map along $(\Id,\Frob):\cM_{d}\to \cM_{d}\times \cM_{d}$ (which is defined since $\cM_{d}$ is smooth equidimensional, see \cite[\S A.1.4]{YZ}) 
%\begin{equation}
%[\Sht^{r}_{\cM_{d}}]:=(\Id,\Frob)^{!}[\Hk^{r, \c}_{\cM_{d}}]\in \Ch_{0}(\Sht^{r}_{\cM_{d}}).
%\end{equation}
%\end{defn}
%
%It seems there is some arbitrariness in extending the fundamental class of $\Hk^{r, \c}_{\cM_{d}}$ to the whole $\Hk^{r}_{\cM_{d}}$, but the next lemma assures that the degree of the resulting $0$-cycle does not depend on the potential choice. 
%
%\begin{lemma}
%Let $\z\in \Ch_{N}(\Hk^{r}_{\cM_{d}})$ be any class that restricts to the fundamental class of $\Hk^{r,\c}_{\cM_{d}}$, and form the Gysin pullback $(\Id,\Frob)^{!}\z\in \Ch_{0}(\Sht^{r}_{\cM_{d}})$. Then for any $(\cE,a)\in \cA_{d}(k)$, we have
%\begin{equation}
%\deg([\Sht^{r}_{\cM_{d}}]|_{\cZ^{r}_{\cE}(a)})=\deg((\Id,\Frob)^{!}\z|_{\cZ^{r}_{\cE}(a)}).
%\end{equation}
%\mnote{Y: may be needed in the comparison with KR cycles.}
%\end{lemma}
%\begin{proof}
%It suffices to show that if $\z'\in \Ch_{N}(\Hk^{r}_{\cM_{d}})$ restricts to zero on $\Hk^{r,\c}_{\cM_{d}}$, then $(\Id,\Frob)^{!}\z'$ has degree zero.  Note that the image of $(\pr_{0},\pr_{r})$ lands in the fiber square $\cM_{d}\times_{\cA_{d}}\cM_{d}$. Consider the Cartesian diagram
%\begin{equation}
%\xymatrix{ \Sht^{r}_{\cM_{d}} \ar[r]\ar[d]^{p} & \Hk^{r}_{\cM_{d}}\ar[d]^{\wt p}\\
%\coprod_{(\cE,a)\in \cA_{d}(k)} \cM_{\cE,a}\ar[r]\ar@{^{(}->}[d] & \cM_{d}\times_{\cA_{d}} \cM_{d}\ar@{^{(}->}[d]\\
%\cM_{d}\ar[r]^-{(\Id,\Frob)} & \cM_{d}\times \cM_{d}
%}
%\end{equation}
%Let $\z''=\wt p_{*}(\z')\in \Ch_{N}(\cM_{d}\times_{\cA_{d}}\cM_{d})$ be the image of $\z'$. By the smallness of $f_{d}$ in Corollary \ref{c:small}, $\cM_{d}\times_{\cA_{d}}\cM_{d}$ has the dimension $N$,  and its top-dimensional components all meet the locus over $\cA_{d}^{\c}$. Since $\z'$ restricts to zero over $\cA^{\c}_{d}$, $\z''=0\in \Ch_{N}(\cM_{d}\times_{\cA_{d}}\cM_{d})$. By base change,  $p_{*}(\Id,\Frob)^{!}\z'=(\Id,\Frob)^{!}\wt p_{*}\z'=(\Id,\Frob)^{!}\z''=0\in \Ch_{0}(\coprod_{(\cE,a)\in \cA_{d}(k)} \cM_{\cE,a})$. Hence $p_{*}((\Id,\Frob)^{!}\z'|_{\cZ^{r}_{\cE}(a)})=0\in\Ch_{0}(\cM_{\cE,a})$. Taking degree (note $\cM_{\cE,a}$ is proper by Corollary \ref{c: hitchin proper})  we conclude that $\deg((\Id,\Frob)^{!}\z'|_{\cZ^{r}_{\cE}(a)})=\deg(p_{*}((\Id,\Frob)^{!}\z'|_{\cZ^{r}_{\cE}(a)}))=0$. 
%\end{proof}
